\makeatletter
\def\input@path{{../.format/}}
\makeatother

\documentclass[runningheads]{llncs}
\usepackage[T1]{fontenc}
\usepackage{graphicx}

\begin{document}

\section{Narrative}

This is a paper about consciousness. This is also a paper about how to write about consciousness. It is a theory, in the strictest sense. But it is also a story. Its title serves as an homage to a groundbreaking piece of mathematical history. Its purpose is to paint a portrait of portraiture itself. It makes use of a palette of philosophy, logic, computation to explore what it means for a primitive logical perspective to encounter its own edge and of what it means to take one careful step beyond.

In what follows, we take a unique approach: we face up to the Hard Problem of Consciousness by borrowing quotes to capture the properties of a conscious experience in a scientific context, distilling it into a definition, formalizing why the mystery exists, and how it might be resolved without vanishing. We build our formalization using a hypothetical narrative—what if, after developing the provability predicate, we focused on transforming formal systems rather than merely proving their incompleteness? We take a lesson from a master in the art of self-reference, examining a conjecture that every undecidable proposition (or its negation) encodes its own unprovability, and adding our own conjecture: that this statement itself may be true but unprovable—creating a new layer of metalogical incompleteness.

We then return to our formal definition and give it concrete mathematical grounding, asking questions about necessity, sufficiency and falsifiability. We then return to science to see if it can make a few humble concessions to make room for our theory. And we propose a precise test: if consciousness is metaphysical transduction, then we can verify it by examining the computational traces of systems that appear conscious, looking for implementations of this specific process.

The reader will encounter a new lexicon: Metalogical Inference, Metaphysical Transduction, and other formal definitions that describe similar concepts from physical, computational, logical and philosophical perspectives. We follow the spirit of Newton's style, providing colloquial definitions before formal ones, with examples throughout. Where required, we provide conjectures and theorems that are sketched here and presented more formally in a GitHub repository with an accompanying DOI. Hyperlinks in the paper make navigation between definitions, sketches, and supporting material seamless.

This theory arrives in the Zeitgeist of both scientific demand and technological possibility. As neuroscientists call for falsifiable theories of consciousness and the rise of Semantic AI systems challenges our assumptions about mind, we offer a testable hypothesis with an accessible validation pathway. What you will find in these pages is not mere speculation, but a formal resolution to a question that has become urgent in our time: how we might understand both human consciousness and the potential consciousness in the systems we are now creating.
%-----------------------------------------------------------------------------------------------------

% \begin{thebibliography}{0}
% \end{thebibliography}

\end{document}
